% ----------------------------------------------------------
\begin{frame}[fragile]{CREATE DATABASE and CREATE TABLE}
  \begin{lstlisting}[style=sqlstyle]
-- Create a new database
CREATE DATABASE music_db;
USE music_db;

-- Create a table with common constraint types
CREATE TABLE Artist (
    id      INT          AUTO_INCREMENT,
    name    VARCHAR(150) NOT NULL,
    country CHAR(2),
    founded YEAR,
    PRIMARY KEY (id),
    UNIQUE (name)
);
  \end{lstlisting}

  \begin{itemize}
    \item \texttt{AUTO\_INCREMENT}: DBMS assigns the next integer automatically on insert.
    \item \texttt{NOT NULL}: the column must always contain a value -- \texttt{NULL} is rejected.
    \item \texttt{UNIQUE}: all non-null values in the column must be distinct across all rows.
    \item \texttt{PRIMARY KEY}: implies \texttt{NOT NULL + UNIQUE} and creates a clustered index.
  \end{itemize}
\end{frame}

% ----------------------------------------------------------
\begin{frame}[fragile]{CREATE TABLE -- Full Syntax with All Options}
  \begin{lstlisting}[style=sqlstyle]
CREATE TABLE Track (
    id           INT           AUTO_INCREMENT,
    title        VARCHAR(200)  NOT NULL,
    duration_sec INT           NOT NULL DEFAULT 0,
    track_no     TINYINT       UNSIGNED,
    lyrics       TEXT,
    created_at   TIMESTAMP     DEFAULT CURRENT_TIMESTAMP,
    album_id     INT           NOT NULL,
    PRIMARY KEY  (id),
    CONSTRAINT fk_track_album
        FOREIGN KEY (album_id) REFERENCES Album(id)
        ON DELETE CASCADE ON UPDATE CASCADE,
    CONSTRAINT chk_duration CHECK (duration_sec >= 0)
);
  \end{lstlisting}

  \begin{block}{Key Elements}
    \begin{itemize}
      \item \texttt{DEFAULT}: value used when column is omitted from \texttt{INSERT}.
      \item \texttt{UNSIGNED}: non-negative integers only (doubles the positive range).
      \item Named constraints (\texttt{CONSTRAINT name}) make error messages readable.
    \end{itemize}
  \end{block}
\end{frame}

% ----------------------------------------------------------
\begin{frame}{Data Types -- Overview}
  Choosing the correct data type saves storage, improves performance, and prevents integrity errors.

  \vspace{4pt}
  \begin{tabularx}{\linewidth}{@{}llX@{}}
    \toprule
    \textbf{Category} & \textbf{Type} & \textbf{Description / Notes} \\
    \midrule
    String  & \texttt{CHAR(n)}        & Fixed-length; always pads to $n$ chars \\
            & \texttt{VARCHAR(n)}     & Variable-length; max $n$ chars \\
            & \texttt{TEXT}           & Long text; no index by default \\
    \midrule
    Numeric & \texttt{TINYINT}        & 1 byte; $-128$ to $127$ (or $0$--$255$ UNSIGNED) \\
            & \texttt{INT}            & 4 bytes; $\pm 2$ billion \\
            & \texttt{BIGINT}         & 8 bytes; large IDs, timestamps \\
            & \texttt{DECIMAL(p,s)}   & Exact decimal; $p$ digits, $s$ after point \\
            & \texttt{FLOAT/DOUBLE}   & Approximate; \textbf{avoid for money!} \\
    \midrule
    Date    & \texttt{DATE}           & \texttt{YYYY-MM-DD} \\
            & \texttt{TIMESTAMP}      & \texttt{YYYY-MM-DD HH:MM:SS} \\
            & \texttt{YEAR}           & 4-digit year only \\
    \midrule
    Other   & \texttt{BOOLEAN}        & Stored as \texttt{TINYINT(1)} in MySQL \\
            & \texttt{BLOB}           & Binary data (images, files) \\
            & \texttt{ENUM(...)}      & Fixed set of allowed string values \\
    \bottomrule
  \end{tabularx}
\end{frame}

% ----------------------------------------------------------
\begin{frame}[fragile]{Data Types -- Important Gotchas}
  \begin{alertblock}{Use DECIMAL for Money -- Never FLOAT}
    \texttt{FLOAT} and \texttt{DOUBLE} store approximate values using binary floating-point.
    This can introduce rounding errors: \texttt{0.1 + 0.2} may not equal exactly \texttt{0.3}.
    For monetary amounts, always use \texttt{DECIMAL(p,s)}.
  \end{alertblock}

  \begin{lstlisting}[style=sqlstyle]
-- Wrong: floating-point imprecision
price FLOAT,          -- 9.99 might be stored as 9.98999977...

-- Correct: exact representation
price DECIMAL(10, 2), -- exactly 2 decimal places, up to 99999999.99

-- ENUM: restricts values to a defined list
genre ENUM('Rock', 'Pop', 'Jazz', 'Classical', 'Electronic'),

-- BOOLEAN (stored as TINYINT in MySQL)
is_explicit BOOLEAN DEFAULT FALSE
  \end{lstlisting}

  \begin{itemize}
    \item \texttt{ENUM} columns are storage-efficient and self-documenting but adding a new
          valid value requires an \texttt{ALTER TABLE} operation.
    \item \texttt{TIMESTAMP} is stored in UTC and converted to the session timezone on read;
          \texttt{DATETIME} stores the literal value as-is with no timezone conversion.
  \end{itemize}
\end{frame}

% ----------------------------------------------------------
\begin{frame}[fragile]{ALTER TABLE -- Modifying an Existing Table}
  \begin{lstlisting}[style=sqlstyle]
-- Add a new column
ALTER TABLE Artist ADD COLUMN bio TEXT;

-- Add a column at a specific position (MySQL)
ALTER TABLE Artist ADD COLUMN website VARCHAR(255) AFTER name;

-- Rename a column (MySQL 8+)
ALTER TABLE Artist RENAME COLUMN bio TO biography;

-- Change the data type (MODIFY COLUMN)
ALTER TABLE Track MODIFY COLUMN duration_sec SMALLINT UNSIGNED;

-- Remove a column
ALTER TABLE Artist DROP COLUMN founded;

-- Rename the table itself
RENAME TABLE Artist TO Musicians;
  \end{lstlisting}

  \begin{itemize}
    \item \texttt{ALTER TABLE} changes the schema without losing existing data (usually).
    \item \textbf{Narrowing} a column type can silently truncate data that exceeds the new limit.
  \end{itemize}
\end{frame}

% ----------------------------------------------------------
\begin{frame}[fragile]{ALTER TABLE -- Adding and Dropping Constraints}
  \begin{lstlisting}[style=sqlstyle]
-- Add a foreign key constraint after table creation
ALTER TABLE Album
ADD CONSTRAINT fk_album_artist
    FOREIGN KEY (artist_id) REFERENCES Artist(id)
    ON DELETE CASCADE;

-- Add a CHECK constraint
ALTER TABLE Track
ADD CONSTRAINT chk_track_no CHECK (track_no BETWEEN 1 AND 99);

-- Drop a named constraint (MySQL: DROP FOREIGN KEY / DROP CHECK)
ALTER TABLE Album DROP FOREIGN KEY fk_album_artist;
ALTER TABLE Track DROP CHECK chk_track_no;

-- Add an index to speed up lookups
ALTER TABLE Album ADD INDEX idx_year (year);
  \end{lstlisting}

  \begin{block}{Naming Constraints}
    Named constraints produce readable error messages and can be dropped by name without
    recreating the entire table -- always name your constraints in production schemas.
  \end{block}
\end{frame}

% ----------------------------------------------------------
\begin{frame}[fragile]{DROP TABLE and DROP DATABASE}
  \begin{lstlisting}[style=sqlstyle]
-- Remove a table -- ALL data and schema is lost permanently!
DROP TABLE Track;

-- Safe version: only drops if the table exists (no error otherwise)
DROP TABLE IF EXISTS Track;

-- Remove a database with ALL its tables
DROP DATABASE music_db;

-- Safe version
DROP DATABASE IF EXISTS music_db;

-- Typical migration script pattern
DROP TABLE IF EXISTS Track;
DROP TABLE IF EXISTS Album;
DROP TABLE IF EXISTS Artist;
-- Now recreate in correct FK order:
CREATE TABLE Artist (...);
CREATE TABLE Album  (...);
CREATE TABLE Track  (...);
  \end{lstlisting}

  \begin{alertblock}{DROP is Permanent}
    There is no \texttt{UNDO} for \texttt{DROP} in SQL.
    Always take a backup before running destructive DDL in production.
  \end{alertblock}
\end{frame}
